\documentclass[12pt]{article}
\usepackage{amsmath, amssymb}
\usepackage{geometry}
\geometry{margin=1in}

\begin{document}

\begin{center}
\textbf{\Large CSE 400 – Fundamentals of Probability in Computing}

\vspace{0.3cm}
\textbf{Scribe: Tutorial-1 Solutions}

\vspace{0.3cm}
School of Engineering and Applied Science (SEAS), Ahmedabad University

Date: February 13, 2026
\end{center}

\hrule
\vspace{0.5cm}

\section*{1. Combinatorics and Multinomial Counting}

\subsection*{Q1. Division of Dishes into Platters}

\subsubsection*{Problem}

20 distinct dishes are to be divided into 4 unlabeled groups of 5 dishes each.

\subsubsection*{(a) Total Number of Ways}

There are 20 distinct dishes divided into 4 groups of 5 dishes each.

Since:
\begin{itemize}
    \item Order within each group does not matter
    \item Order of the groups does not matter
\end{itemize}

Using the multinomial formula:

\[
\text{Total ways} =
\frac{20!}{(5!)^4 4!}
\]

\subsubsection*{(b) Probability All Platters Have Same Cuisine}

Each cuisine contains exactly 5 dishes.

The only way each platter can contain dishes from a single cuisine is if each cuisine forms its own platter.

Because platters are unlabeled, this arrangement is unique.

\[
P =
\frac{1}{\frac{20!}{(5!)^4 4!}}
=
\frac{(5!)^4 4!}{20!}
\]

\section*{2. Uniform Distribution and Geometric Probability}

\subsection*{Q2. Bus Waiting Time}

Arrival time is uniformly distributed over a 30-minute interval.

\subsubsection*{(a) Waiting Time $< 5$ Minutes}

Favorable intervals:
\begin{itemize}
    \item 7:10–7:15
    \item 7:25–7:30
\end{itemize}

Total favorable time = 10 minutes.

\[
P(\text{wait} < 5) = \frac{10}{30} = \frac{1}{3}
\]

\subsubsection*{(b) Waiting Time $> 10$ Minutes}

Favorable intervals:
\begin{itemize}
    \item 7:00–7:05
    \item 7:15–7:20
\end{itemize}

Total favorable time = 10 minutes.

\[
P(\text{wait} > 10) = \frac{10}{30} = \frac{1}{3}
\]

\section*{3. Conditional Probability and De Morgan’s Law}

\subsection*{Q3. Late Arrival and Early Leaving}

Let:
\begin{itemize}
    \item $L$ = arrives late
    \item $E$ = leaves early
\end{itemize}

Given:

\[
P(L) = 0.15,\quad P(E) = 0.25,\quad P(L \cap E) = 0.08
\]

Arriving early = $L^c$

Leaving late = $E^c$

\[
P(E^c) = 1 - 0.25 = 0.75
\]

\[
P(L \cup E) = 0.15 + 0.25 - 0.08 = 0.32
\]

Using De Morgan’s Law:

\[
P(L^c \cap E^c) = 1 - 0.32 = 0.68
\]

\[
P(L^c \mid E^c) =
\frac{0.68}{0.75}
\approx 0.907
\]

\section*{4. Poisson Distribution}

\subsection*{Definition}

If $X \sim \text{Poisson}(\lambda)$, then:

\[
P(X = k) =
\frac{e^{-\lambda} \lambda^k}{k!}
\]

\subsection*{Q4. Typographical Errors}

\[
\lambda = 0.2
\]

\subsubsection*{(a) $P(X=0)$}

\[
P(X=0) = e^{-0.2} = 0.8187
\]

\subsubsection*{(b) $P(X \ge 2)$}

\[
P(X \ge 2) = 1 - P(0) - P(1)
\]

\[
P(1) = 0.2e^{-0.2} = 0.1637
\]

\[
P(X \ge 2) = 1 - 0.8187 - 0.1637
\]

\[
P(X \ge 2) = 0.0176 \approx 0.0175
\]

\subsection*{Q5. Airplane Crashes}

\[
X \sim \text{Poisson}(3.5)
\]

\[
P(X=0) = 0.0302
\]

\[
P(X=1) = 0.1057
\]

\subsubsection*{(a) $P(X \ge 2)$}

\[
P(X \ge 2) = 1 - 0.0302 - 0.1057
\]

\[
P(X \ge 2) = 0.8641
\]

\subsubsection*{(b) $P(X \le 1)$}

\[
P(X \le 1) = 0.0302 + 0.1057 = 0.1359
\]

\subsection*{Q6. Highway Accidents}

\[
X \sim \text{Poisson}(3)
\]

\[
P(0) = 0.0498
\]

\[
P(1) = 0.1494
\]

\[
P(2) = 0.2240
\]

\subsubsection*{(a) $P(X \ge 3)$}

\[
P(X \ge 3) = 1 - 0.4232 = 0.5768
\]

\subsubsection*{(b) Conditional Probability}

\[
P(X \ge 1) = 0.9502
\]

\[
P(X \ge 3 \mid X \ge 1)
=
\frac{0.5768}{0.9502}
\approx 0.607
\]

\section*{5. Gaussian Distribution and Q-Function}

\subsection*{Q7. Gaussian Random Variable}

Given:

\[
f_X(x) =
\frac{1}{\sqrt{8\pi}}
\exp\left(
-\frac{(x+3)^2}{8}
\right)
\]

Comparing with Gaussian PDF:

\[
m = -3,\quad \sigma^2 = 4,\quad \sigma = 2
\]

Gaussian relations:

\[
F_X(x) = \Phi\left(\frac{x-m}{\sigma}\right)
\]

\[
P(X > x) = Q\left(\frac{x-m}{\sigma}\right)
\]

\begin{enumerate}
\item $P(X \le 0)$

\[
\Phi\left(\frac{3}{2}\right)
=
1 - Q\left(\frac{3}{2}\right)
\]

\item $P(X > 4)$

\[
Q\left(\frac{7}{2}\right)
\]

\item $P(|X+3|<2)$

\[
-5 < X < -1
\]

\[
= \Phi(1) - \Phi(-1)
\]

\[
= 2\Phi(1) - 1
= 1 - 2Q(1)
\]

\item $P(|X-2|>1)$

\[
X < 1 \text{ or } X > 3
\]

\[
= 1 - Q(2) + Q(3)
\]

\end{enumerate}

\section*{6. Binomial Distribution}

\subsection*{Q8. Jury Decision}

If defendant is innocent:

\[
\sum_{i=5}^{12}
\binom{12}{i}
\theta^i (1-\theta)^{12-i}
\]

If defendant is guilty:

\[
\sum_{i=8}^{12}
\binom{12}{i}
\theta^i (1-\theta)^{12-i}
\]

Let $\alpha$ be probability defendant is guilty.

\[
\alpha
\sum_{i=8}^{12}
\binom{12}{i}
\theta^i (1-\theta)^{12-i}
+
(1-\alpha)
\sum_{i=5}^{12}
\binom{12}{i}
\theta^i (1-\theta)^{12-i}
\]

\section*{7. Discrete Random Variables}

\subsection*{Q9. Given}

\[
p(i) = c\frac{\lambda^i}{i!}
\]

Since

\[
\sum_{i=0}^{\infty} p(i) = 1
\]

\[
c e^\lambda = 1
\]

\[
c = e^{-\lambda}
\]

\subsubsection*{(a)}

\[
P(X=0) = e^{-\lambda}
\]

\subsubsection*{(b)}

\[
P(X>2)
=
1 - e^{-\lambda} - \lambda e^{-\lambda} - \frac{\lambda^2 e^{-\lambda}}{2}
\]

\subsection*{Cumulative Distribution Function}

\[
F(a) = \sum_{x \le a} p(x)
\]

If $X$ takes values $x_1 < x_2 < x_3 < \dots$, then $F$ is a step function with jumps of size $p(x_i)$.

Example provided:

\[
p(1)=\frac14,\;
p(2)=\frac12,\;
p(3)=\frac18,\;
p(4)=\frac18
\]

Corresponding stepwise CDF as specified.

\section*{8. System Reliability}

\subsection*{Q10. Majority System}

Number functioning is binomial $(n,p)$.

\subsubsection*{(a) 5 vs 3 Components}

5-component effective:

\[
10p^3(1-p)^2 + 5p^4(1-p) + p^5
\]

3-component effective:

\[
3p^2(1-p) + p^3
\]

Condition reduces to:

\[
3(p-1)^2(2p-1) > 0
\]

\[
p > \frac12
\]

\subsubsection*{(b) General Case}

\[
P_{2k+1} - P_{2k-1}
=
\binom{2k-1}{k}
p^k(1-p)^k (2p-1)
\]

Since binomial coefficient and powers are positive,

\[
> 0 \iff p > \frac12
\]

\vspace{0.5cm}
\begin{center}
\textbf{End of Scribe}
\end{center}

\end{document}